\chapter{はじめに}

\section{目的}

本指導書の目的は次の通りである。

\begin{itemize}
    \item 情報技術として広く使われている人工ニューラルネットワークをプログラム実装を通じて理解すること
    \item 基本的なアナログ部品であるオペアンプを用いたアナログ回路について、基本的なふるまいをシミュレーションおよび電子回路の計測を通じて理解すること
    \item 人工ニューラルネットワークの基本である多層パーセプトロンによる識別器を実装することで、ソフトウェアおよびハードウェアを協調した設計について理解すること
    \item 多人数による開発設計を通じて技術的なコミュニケーション、およびプロジェクト進行について理解すること
\end{itemize}

\section{本実験の流れ}

合計6週で実施する。
大まかなスケジュール、および提出物について表 \ref{table:schedule} に示す。

\begin{table}[htbp]
  \centering
  \caption{実験スケジュールおよび提出物一覧}
  \label{tab:schedule}
  \vspace{0.5em}
  % 全体の横幅を本文領域（\textwidth）に合わせ、各列の比率を調整
  \begin{tabular}{|c||c|c|c|}
    \toprule
    \textbf{回数} & \textbf{実施内容} & \textbf{分類} & \textbf{提出物} \\
    \midrule
    1 & 人工ニューラルネットワークの実装 & 個人 & 小テスト（計算問題） \\
    2 & 逆誤差伝搬法による学習の実装 & 個人 & レポート\newline（ソフトウェア実装と解析） \\
    3 & オペアンプによるMLPの実装 & 個人 \& ペア & レポート\newline（回路実装と評価） \\
    4 & MNIST識別器の設計（SW\&HW） & グループ & --- \\
    5 & MNIST識別機の設計（SW\&HW） & グループ & --- \\
    6 & プレゼンテーション & グループ & レポート \\
    \bottomrule
  \end{tabular}
\end{table}